\documentclass[12pt]{article}

\usepackage{amsmath}
\usepackage{amssymb}
\usepackage{amsthm}
\usepackage{graphicx}
\usepackage{float}

\title{CGA Math}
\author{Spencer T. Parkin}

\newcommand{\nvao}{o}
\newcommand{\nvai}{\infty}

\begin{document}
\maketitle

This paper documents the math needed to compose and decompose elements of CGA.

\section*{Foundation}

For any blade $B$ of CGA, it represents two geometries.  These are all vectors $x$ in the euclidean sub-algebra of CGA
such that
\begin{equation*}
p(x)\cdot B = 0,
\end{equation*}
or all such vectors with
\begin{equation*}
p(x)\wedge B = 0,
\end{equation*}
where $p(x)$ is defined as the vector
\begin{equation*}
p(x) = \nvao + x + \frac{1}{2}x^2\nvai.
\end{equation*}
We'll call the first of these the ``dot representation'' and the second, ``wedge representation''.
To convert from one to the other, you can multiply by $I$, the unit psuedo-scalar.
This is because
\begin{equation*}
p(x)\cdot B = 0\iff p(x)\wedge BI=0.
\end{equation*}
Of course, $B$ and $BI$ are duals of one another, so the conversion goes both ways.

Using blades to represent geometries, we get the following useful properties.
If $B_1$ and $B_2$ are blades with $B_1\wedge B_2\neq 0$, then we have
\begin{equation*}
\{x|p(x)\cdot B_1\wedge B_2 = 0\} = \{x|p(x)\cdot B_1=0\}\cap\{x|p(x)\cdot B_2=0\},
\end{equation*}
as well as
\begin{equation*}
\{x|p(x)\wedge B_1\wedge B_2 = 0\} \supseteq \{x|p(x)\wedge B_1=0\}\cup\{x|p(x)\wedge B_2=0\}.
\end{equation*}
In this latter case, since there are only so many shapes in CGA, we can reasonably deduce that
$B_1\wedge B_2$ will be the geometry that fits those represented by $B_1$ and $B_2$.  You may,
for example, find the sphere that fits a given circle and point, or the circle that fits 3 given points,
and so on.

The shapes of CGA are either round or flat.  If using the wedge presentation, a round can be converted
into a flat by tacking on $\nvai$ with the outer product.  For example, if $B$ represents a
circle using the wedge representation, then $B\wedge\nvai$ is the plane containing the circle,
also in the wedge representation.

\section*{Compositions}

For all geometry that follows, we let $c$ and $n$ denote euclidean vectors, and $r$ denote a scalar.
We always require that $|n|=1$.

We are going to use the dot representation here instead of the wedge representation,
and ignore the magnitude of $B$ as long as it is non-zero.

\subsection*{Planes}

Letting $B=n + (n\cdot c)\nvai$, we have
\begin{align*}
0 &= p(x)\cdot B \\
 &= \left(\nvao + x + \frac{1}{2}x^2\nvai\right)\cdot(n + (c\cdot n)\nvai) \\
 &= -c\cdot n + x\cdot n \\
 &= (x-c)\cdot n,
\end{align*}
which is clearly the equation for a plane centered at $c$ with normal $n$.

\subsection*{Lines}

Lines are a bit trickier.  Here, let $i=e_1e_2e_3$ denote the unit psuedo-scalar of the 3D euclidean sub-algebra of CGA.
Doing so, let $B=ni + (ni\cdot c)\nvai$.  We then see that
\begin{align*}
0 &= p(x)\cdot B \\
 &= \left(\nvao + x + \frac{1}{2}x^2\nvai\right)\cdot (n - (c\cdot ni)\nvai) \\
 &= \left(\nvao + x + \frac{1}{2}x^2\nvai\right)\cdot (n + \nvai\wedge (c\cdot ni)) \\
 &= (\nvao\cdot\nvai)(c\cdot ni) + x\cdot ni - (x\cdot(c\cdot ni))\nvai \\
 &= -c\cdot ni + x\cdot ni - (x\cdot (c\wedge n)i)\nvai \\
 &= (x - c)\cdot ni - (x\wedge c\wedge n)i\nvai.
\end{align*}
Here, we must have $0=(x-c)\cdot ni$ as well as $0=x\wedge c\wedge n$.
The first of these is clearly a line centered at $c$ with normal $n$, which is what we're after, but what of the second?
Well, we need only consider those values of $x$ for which the first is satisfied.
If $x=c$, we have a trivial case, and all is well.
Otherwise, we have $x-c=\lambda n$ for some non-zero scalar $\lambda$,
which means that $x\wedge c\wedge n=x\wedge c\wedge\frac{1}{\lambda}(x-c)$, which
is clearly zero.

\subsection*{Spheres}

Consider $B=\left(\nvao + c + \frac{1}{2}(c^2 - r^2)\nvai\right)$.  We then have
\begin{align*}
0 &= p(x)\cdot B \\
 &= \left(\nvao + x + \frac{1}{2}x^2\nvai\right)\cdot\left(\nvao + c + \frac{1}{2}(c^2 - r^2)\nvai\right) \\
 &= -\frac{1}{2}(c^2 - r^2) + x\cdot c - \frac{1}{2}x^2.
\end{align*}
Rearranging this, we find that
\begin{align*}
r^2 &= c^2 - 2(x\cdot c) + x^2 \\
 &= (x-c)^2,
\end{align*}
which is clearly the equation for a sphere.

Interestingly, we also find imaginary spheres in CGA with $B=\left(\nvao + c + \frac{1}{2}(c^2 + r^2)\nvai\right)$.
Imaginary rounds are necessary in cases where an intersection is computed that can't actually exist.

\subsection*{Points}

Points are simply the special case of spheres with $r=0$.  A point at $c$ is simply $p(c)$.
Note that these, technically, are round points, not to be confused with tangent points
or flat points.

For a flat point, let
\begin{equation*}
B = (n + (n\cdot c)\nvai)\wedge(ni + (ni\cdot c)\nvai),
\end{equation*}
which is clearly the intersection of a line and a plane.  If you were to expand
this wedge product out, things would cancel, and you'd be left with
\begin{equation*}
B = i + ci\wedge\nvai,
\end{equation*}
showing that indeed $n$ is not applicable here.

\subsection*{Circles}

Here, let
\begin{equation*}
B=\left(n + (n\cdot c)\nvai\right)\wedge\left(\nvao + c + \frac{1}{2}(c^2-r^2)\nvai\right),
\end{equation*}
which is the intersection of a sphere centered on a plane.  Use an imaginary sphere here
to get an imaginary circle.  Such circles form the intersection of, for example, two
real spheres that are intersected, but that don't actually intersect.

\subsection*{Point-Pairs}

Point-pairs may seem odd, but they're simply the 1-dimensional analog of spheres and circles.
Here, let
\begin{equation*}
B=(ni+(ni\cdot c)\nvai)\wedge\left(\nvao + c + \frac{1}{2}(c^2 - r^2)\nvai\right).
\end{equation*}
This is simply the intersection of a line going through the center of a sphere.
Similar to cricles, use an imaginary sphere here to get an imaginary point-pair.
Such a point-pair would occur, for example, when you compute the intersection
of a line and a sphere that don't actually intersect.

\subsection*{Tangent Points}

These are the special case of circles or point-pairs when $r=0$.
Note that while it takes 4 points to fit a sphere, it only takes 2 tangent points, or 2 points and 1 tangent point, to do the same.
You can get a tangent point by intersecting two spheres that touch in a single point.

\section*{Decomposition}

Here we utilize the composition forumlas derived earlier in order to extract the variables $c$, $n$ and $r$
from a given blade $B$.  Since we can't ignore magnitude here, we need an extra scalar $w$
to represent the ``weight" of the geometry.

\section*{Transformations}

\subsection*{Rotors}

\subsection*{Motors}

\subsection*{Transversions}

\end{document}